% \setchapterpreamble[u]{\margintoc}
\chapter{Introduction}
\labch{intro}


% In the next few years, the internet will fill up with AI generated slop. To be fair, much of the internet is already slop anyways... 
In the library of Babel\sidenote{\textit{La biblioteca de Babel}, is a short story by Argentine author Jorge Luis Borges, inspired by German author Kurd Lasswitz in his 1901 story \textit{Die Universalbibliothek}. Interestingly, technology has enabled this fictional construct to virtually materialize: head to \url{https://libraryofbabel.info/} to browse it.}, the books collectively contain the entirety of human knowledge. Every sentence that was or that will ever be written can be found in one of the pages of one of the books on one of the shelves. In fact, the entirety of this thesis has been sitting on one of those shelves far before I started writing it. Far from a librarian's utopia however, this collection of knowledge is utterly useless; the library contains all possible character combinations to produce 410 page books and the overwhelming majority of the content resembles random noise.\sidenote{Attempts at compressing this knowledge have sadly still been unsuccessful: see \autoref{fig:hbs_meme}} It matters not having access to all knowledge, what matters is being able to tractably find the \textit{bits} that matter.


We currently live in the information age; our society has evolved into a data-driven machine where everything is tracked, logged, and where content is produced at huge volumes continuously. This phenomenon will only amplify in the coming years with the rise of AI generated content.
As the price of data storage reaches new lows every year\sidenote{Kryder's law prescribes that the amount of storage space available is growing approximately exponentially.}, the incentives to curate, modify or delete existing content are often overwhelmed by the ease of storing irrationally, and the facility of producing new content - often devoid of meaningful new signal. The internet is not exactly a library of Babel, but it tends towards it.

Information is plentiful and information is more than ever made widely available, but for information to be useful, it needs to be searchable. As volume increases, separating \emph{signal} (relevant information) from \emph{noise} is crucial. The field of Information Retrieval is about developing modern means to thread \emph{"through the consequent maze to the momentarily important item"}.\sidenote{This sentence is a direct reference to an essay, \emph{As We May Think} (1945) \cite{bush1945aswemaythink}, in which Vannevar Bush describes the state of information retrieval at the time: "The summation of human experience us being expanded at a prodigious rate, and the means we use for threading through the consequent maze to the momentarily important item is the same as was used in the days of square-rigged ships."} 

% \textit{The summation of human experience is being expanded at a prodigious rate}, but \textit{the means we use for threading through the consequent maze to the momentarily important item}\sidenote{This sentence is a direct reference to an essay, \emph{As We May Think} (1945) \cite{bush1945aswemaythink}, in which Vannevar Bush describes the state of information retrieval at the time: "The summation of human experience us being expanded at a prodigious rate, and the means we use for threading through the consequent maze to the momentarily important item is the same as was used in the days of square-rigged ships."} are of crucial importance. The words of Vannevar Bush echo information theory: knowledge transmission requires separating signal from noise. 



Over time, tools have been developed to make sense of the information deluge. Around 240 BC, Callimachus of Cyrene, a Greek poet, attempted to catalog the library of Alexandria in his now lost \textit{Pinakes}, the first knowledge index known to man. Millennia later, the largest information repository is now the internet, and the even more daunting task of its indexation has led to the rise and fall of certain modern day empires: large corporations\sidenote{Over 5 trillion Google searches were made in 2025. \cite{blogPersonalizationFuture}}. In between both of these herculean enterprises, the advent of computing has been the largest catalyst of progress to enable efficiently searching through large quantities of data.


Vannevar Bush\sidenote{Vannevar Bush (1890–1974) was an American engineer and science administrator who led U.S. wartime R\&D through the U.S. Office of Scientific Research and Development. As World War 2 ended, he turned to a peacetime problem: how science could keep advancing amid an exploding record of knowledge. In “As We May Think” (1945), his proposed memex anticipated associative information retrieval.}, a visionnary scientist who headed  the U.S. Office of Scientific Research and Development during the second world war, describes \emph{search} at a time before widespread computing:

\begin{lavenderkaobox}
    "When data of any sort are placed in storage, they are filed alphabetically or numerically, and information is found (when it is) by tracing it down from subclass to subclass. [...] The human mind does not work that way. It operates by association."
    
    \flushright -- Vannevar Bush, As We May Think (1945)
\end{lavenderkaobox}

Crucially, computing has since enabled transitioning away from fixed document classification systems designed to accommodate all future searches (catalogs) and allowed for the creation of more dynamic systems that adapt to individual queries (eg. keyword search systems). True to Bush's vision, such systems enable richer associative linking between queries and documents, and were only made possible due to technological advances. Later demonstrated in the 1960s by the Cranfield experiments \sidecite{cleverdon1967cranfield} \sidenote{The Cranfield experiments were a series of experimental studies in information retrieval conducted by Cyril W. Cleverdon at the College of Aeronautics, today known as Cranfield University, in the 1960s to evaluate the efficiency of indexing systems.}, these word frequency based systems both yield better results and are drastically easier to implement.


This thesis focuses on an even more recent development of modern search systems: the use of deep neural networks applied to \textit{search} and document retrieval. Leveraging semantic associative links, neural search is today our best tool to enable humans (and AIs) to augment their capacities by leveraging external memories. 




\section{This Thesis in a Glance} 

Information retrieval (IR)\sidenote{IR is used interchangeably with \emph{Search} in this thesis.} \emph{is the process of finding material (usually documents) of an unstructured nature (usually text) that satisfies an information need from within large collections (nowadays, usually stored on computers)}.\sidenote{This definition is largely inspired by \cite{manning2008introduction_ir}.} In simpler terms, it is the process of using a query to find and rank information resources (broadly called \emph{documents}\sidenote{A \textbf{document} is an information unit of limited size, typically small enough for a human reader to assess its relevance to a query, such as a page, paragraph, or passage. Although documents are ideally self-contained, they often depend on surrounding context; for example, a page in a textbook may contain the sought-after formula while requiring preceding pages for proper interpretation.}) that are relevant to a user’s information need.

Today, the deep neural networks used to index and query information mostly rely on a \emph{semantic understanding} of the content to retrieve documents at large scales. This contrasts with previous methods that leveraged meta-information (title, author, tags, document visits), or relied on statistical heuristics to match content to a user query (keyword matching). 

More precisely, efforts in deep semantic search have to this day mostly focused on creating meaningful latent representation spaces to map short text passages to their equivalent latent representations (\emph{embeddings}) in a manifold that aligns semantic proximity with geometric distance. This thesis argues that in practice, solely viewing AI search through this prism is quite restrictive: search systems should move beyond standard text embeddings to perform document retrieval. The thesis explores two main improvement directions which share one common principle: leveraging easily available yet often overlooked information (\emph{priors}) to strengthen associative links between queries and relevant documents. An important design principle is that the technical propositions aim to reduce inductive bias and be compatible with model and data scaling such that they should remain relevant with respect to future improvements in the field \sidecite{sutton2019bitter}. 


% This thesis argues AI-powered search systems should move beyond standard text embeddings to perform document retrieval, and proposes several improvement directions, which share a common principle of  leveraging easily available yet often overlooked information (priors) that can aid information retrieval objectives. 





\subsection{Part 1: Visual Document Retrieval} 
% \subsection{Part 1: Retrieving Documents Requires Vision}


Historically, AI search engines have treated documents primarily as text. This thesis introduces the concept of \textit{Visual Document Retrieval} (VDR) through the publication ``ColPali: Efficient Document Retrieval with Vision Language Models'' \sidecite{faysse2025colpaliefficientdocumentretrieval}. This work proposes performing document retrieval without relying on text extracted from document pages. Instead, pages are indexed directly from their images, allowing the retrieval system to capture the semantic content expressed through pixels, including text, layout, tables, figures, and other visual elements.

We show that this approach substantially simplifies and accelerates document indexing pipelines, while also improving retrieval performance, particularly for visually complex documents with rich layouts, tables, figures, or other non-linear structures (Chapter \autoref{ch:colpali}).

The success of VDR has led since the original publication to many follow up works by academic and industrial actors. As a result, model performances rapidly improved, saturating the original Visual Document Retrieval benchmark \textit{ViDoRe} introduced in the ColPali paper, and leading to subsequent benchmarking efforts. In ViDoRe V2 \sidecite{macé2025vidorebenchmarkv2raising}, we construct complex and realistic pairs of queries and documents to assess retrieval performance on open-ended questions spanning up to multiple pages and rely on complex AI pipelines to simulate realistic interactions between users and information corpora. In ViDoRe V3 \sidecite{loison2026vidorev3comprehensiveevaluation}, we collaborate with Nvidia, to further scale annotation efforts on varied domains reflecting real world use of industrial search systems and develop an annotation pipeline that combines AI agents with over 12,000 hours of expert human annotators. These contributions are detailed in Chapter \autoref{ch:vidorev3}.

A year after the publication of the ColPali paper, we released \emph{ModernVBert}, a visual document retrieval model that commoditizes VDR performance by showcasing similar performances than the original ColPali model while being 10x smaller and running on CPU \sidecite{teiletche2025modernvbert}. This was made possible by through ablations of all performance factors behind vision language retrieval models and many critical insights were uncovered and detailed in Chapter \autoref{ch:modernvbert}.


\subsection{Part 2: Contextual Document Retrieval}

Another limitation of standard embedding-based search systems is that queries and documents are usually processed and embedded in isolation, without relying on external context which is in practice usually available and often crucial. 

\noindent\textbf{Document Retrieval Requires Context.} Book pages are not meant to be read in isolation. In practice however, this is what standard indexing pipelines do: they split long documents into independent chunks and embed them individually. 
This shortcoming is visible in current visual retrievers, but is also a major weakness of the more common textual embedding models which embed long documents as a series of independent passages, without exploiting the fact to a reader, certain passages might be key to understand others. In Chapter \autoref{ch:conteb}, we quantify the impact of context aware retrieval and propose technical solutions based on late-chunking aware training to infuse contextual information into chunk embeddings \sidecite{conti2025context}. While this work is instantiated on textual retrievers for ease of experimentation, it is directly applicable to our work on improving visual retrievers.


\noindent\textbf{Query Understanding Requires Context.}
The lack of contextualization does not lie solely on the document side. To be properly parsed, queries also require context (what does the user really look for), and often require multiple steps of reasoning and background knowledge to be properly processed. As a simple example, the query \emph{"What is the length of the longest French river ?"}, you first need to know the river is the Loire, then can proceed to search for the Loire's length. Standard embedding based search systems would be mostly incapable of such multi-step reasoning, hence we turn to the most recent evolution of search engines: AI agents powered by Large Language Models (LLMs). Roughly defined, these are systems in which a language model acts as an orchestrator, decomposing a user request into intermediate steps, issuing retrieval or tool-use actions, and integrating the results into a final answer.

Agents leverage knowledge (priors) encoded in their parametric knowledge ("weights") to contextualize queries and documents with previously memorized information. In a sense, query terms trigger associative recall of internal knowledge through the forward pass mechanism. Studying the mechanisms behind LLM memorization is key to understand how domain specific pre-training can impact parametric memorization and downstream performance. Through training, we can increase the chance useful \emph{priors} will be encoded in the model weights and help contextualize queries at inference time.

% Recent efforts such as BrowseComp highlight the growing importance of AI agents for real-world information seeking, and suggest that robust access to practically useful information increasingly benefits from \emph{adaptive compute}. Importantly, these systems continue to depend on strong first-layer retrieval primitives---notably embedding-based search---to identify promising candidates before deeper reasoning, verification, or tool use can take place, and should be seen as a complementary abstraction.


\noindent\textbf{Improving LLMs is Improving Search.} Agentic search largely benefits from capable LLM orchestrators. More capable models are infused with better reasoning skills and better background knowledge enabling a more precise use of retrieval tools and are able to retrieve information quicker and in less turns (BrowseCompPlus). While the best LLMs have historically been proprietary models of large size, essentially pretrained on the English language, we argue smaller models trained openly and transparently are capable of strong performance and benefit the academic and industrial community. In Chapter \autoref{ch:croissant}, we introduce CroissantLLM \sidecite{faysse_croissantllm_2024}, a compact bilingual French--English language model trained on an equal mixture of English and French data, achieving state of the art performance for LLMs of the size at the time of release.


% Finally, a line of work conducted as part of this thesis addresses the transparency problem from a more adversarial perspective. Led primarily by the Computational Privacy Group at Imperial College London, the two papers are discussed only briefly in Chapter~X. They focus on \emph{training-data memorization} and \emph{membership inference} in large language models: to what extent are individual training examples memorized, and how can one provide evidence that a given text was included in an LLM's training corpus? These questions have direct relevance to ongoing legal disputes and raise broad policy implications, notably around copyright and data governance. In \sidecite{meeus2024copyright}, we propose inserting a small fraction of controlled ``canaries'' into the training data in order to reliably characterize the regimes in which memorization occurs and to quantify the associated risk of information leakage. This work is done through the CroissantLLM training project. Complementarily, \sidecite{meeus2024sok} provides a systematic  overview of membership inference attacks against LLM training data, organizing existing methods, clarifying threat models and assumptions, and highlighting open problems. Together, these contributions move toward more rigorous and auditable assessments of modern LLM behavior, and offer practical tools to reason about privacy and provenance in large-scale training pipelines.


\subsection{Conclusion: Perspectives On Search}

What does the next generation of search systems look like? Is Retrieval Augmented Generation (RAG) a thing of the past, buried by the newest shiny model or technique?



\begin{marginfigure}
    \centering
    \includegraphics{images/misc/rag.png}
    \caption[RAG is dead meme]{Retrieval Augmented Generation (RAG) has been buried a dozen times according to LinkedIn AI influencers, replaced by agents, long context models, keyword search... In practice, models will change, but the need to condition a generative process on external information will remain very much alive.}

    % Latencies indicated in the figure are computed on the same machine \texttt{g2-standard-16} with a NVIDIA-L4 GPU, and averaged over 100 pages. Further details in \autoref{subsection:latency}.The computed score is the average of from all datasets in \textit{ViDoRe} for PDF Parser + Captioning and ColPali respectively. Further details in \autoref{section:results}.} 
    \label{fig:rag_dead}
\end{marginfigure}



At its core, indexing and retrieving information from external sources in a working "memory" is likely to remain one of the key challenges of AI systems in the future. Scientific advances will unlock better recall from broader sources and more personalized answers and the form factor will change, but the need to rely on external verifiable non-parametric resources to supplement our memory, and the need to \emph{navigate the maze} and separate information from noise will remain key.  

The parting words in Chapter \autoref{ch:conclusion} offer some personal perspectives into what I believe to be impactful future research directions in the retrieval field. In line with the rest of this thesis, I preach efforts into omni-modal universal embeddings, scaffolding engineering, reflect on ways to push benchmarking efforts forward, and detail why I believe bridging efforts with other NLP fields such as long context modeling research might enable scaling the training of embedding models multiple orders of compute magnitude. To anyone reading this thesis and wishing to discuss these ideas with me, my mail is open.


\section{List of Publications and Impact} 

This thesis is recollection of a focused selection of first-author\sidenote{I am last author on Vidore V2 and ModernVBert but was the one setting the research direction, closely supervising the projects and mostly writing the produced manuscripts} papers produced during the first 2.5 years of my PhD program. Most papers presented are inserted verbatim as chapters of the thesis and left unmodified with respects to the original peer-reviewed publication of the work. 


\noindent\textbf{Visual Document Retrieval.} The first part focuses on visual document retrieval, beginning with 
\emph{ColPali: Efficient Document Retrieval with Vision Language Models} 
\sidecite{faysse2025colpaliefficientdocumentretrieval}, published at ICLR 2025. It then turns to 
\emph{ModernVBERT: Towards Smaller Visual Document Retrievers} 
\sidecite{teiletche2025modernvbert}, published at ICML 2026, and to the ViDoRe benchmark series: 
\emph{ViDoRe Benchmark V2: Raising the Bar for Visual Retrieval} and 
\emph{ViDoRe V3: A Comprehensive Evaluation of Retrieval Augmented Generation in Complex Real-World Scenarios} 
\sidecite{macé2025vidorebenchmarkv2raising, loison2026vidorev3comprehensiveevaluation}, accepted at ACL 2026\sidenote{Although I initiated the project and was the original senior author, my role became more advisory towards the end as I joined the team at Meta FAIR working on unrelated topics for the last 6 months of my PhD.}. 


The contributions of this thesis go beyond a series of peer reviewed publications. Significant efforts were devoted to continuously improve the performance, to extend the capabilities and to advocate for the use of Visual Document Retrieval models. We can notably mention the release of the ColQwen model series\sidenote{We published the models openly and integrated the code on HuggingFace \url{https://huggingface.co/docs/transformers/main/model_doc/colqwen2}}, a series of SOTA visual retrievers, downloaded over 5 million times on HuggingFace at the time of writing. We also extended retrieval to other modalities such as video and audio through ColQwen-Omni \sidenote{We published the model openly on the HuggingFace platform and released a blog post at \url{https://huggingface.co/blog/manu/colqwen-omni-omnimodal-retrieval}.} and maintained the ColPali GitHub library at \url{https://github.com/illuin-tech/colpali}. The field of Visual Document Retrieval is now well established and thriving.



\noindent\textbf{Contextual Document Retrieval.} The second part focuses on leveraging contextual information to improve retrieval. It covers 
\emph{Context is Gold to find the Gold Passage: Evaluating and Training Contextual Document Embeddings} 
\sidecite{conti2025context}, accepted as an Oral at EMNLP 2025, and 
\emph{CroissantLLM: A Truly Bilingual French-English Language Model} 
\sidecite{faysse_croissantllm_2024}, published at TMLR 2025.

A FEW IMPACT METRICS ?

% exclude from thesis 

\noindent\textbf{Other First Author Papers.} Thesis manuscripts are focused by nature. My thesis itself was less so. Towards the end, a small section is dedicated to presenting a few papers linked to this main body of work. Notably, in "Self-Pruned Key Value Attention: Learning When To Write by Predicting Future Utility" \sidecite{szilvasy2026selfprunedkeyvalueattentionlearning} currently undergoing peer review, we tackle in-context retrieval problems through long context language modeling. Viewing a LLM KV Cache as a memory mechanism, we show we can induce the LLM to learn which information to store or discard, leading to reduced memory usage and improved speeds with no performance degradations. In another first author paper, "Revisiting Instruction  Model Evaluation" \sidecite{faysse2023revisiting} accepted as an Oral at EMNLP 2023, we show the benefits of using LLMs as judges for capability assessment. While this approach is now widespread\sidenote{Written in May 2023 and released in October 2023 after the EMNLP anonymity period, following the rules in place at the time.}, this paper was an early work advocating the use of LLM-based evaluation for broad capability measurement as they display strong task agnostic grading scales.
%The paper stems from the observation that instruction-tuned LLMs are deployed across an increasingly diverse set of tasks, making task-specific metrics difficult to maintain, poorly calibrated, and ultimately non-scalable. It shows that LLM judges can provide consistent signal across heterogeneous tasks and remain robust to variation in expected answer formats---settings where many traditional generative metrics fall short. The analysis further reveals distinct phases during supervised fine-tuning: models often learn to match the \emph{format} of the desired output before improving underlying task performance. Building on this finding, the work proposes leveraging synthetic data augmentation to improve data efficiency. Overall, this chapter argues for a fairer and more practical assessment of the capabilities of this new generation of models. Since the release of this work, LLM-as-a-judge methods have become a dominant approach for evaluating RAG systems and related techniques \sidecite{es2025ragasautomatedevaluationretrieval, kim2023prometheus, muller2025grousebenchmarkevaluateevaluators}.


\noindent\textbf{Other Papers.}  During this thesis, I also contributed as an active but secondary author to many other publications. Notably, in \sidecite{gisserotboukhlef2024trustworthyrerankingsimpleeffective} (TMLR 2024) we study whether the distribution of retrieval score can be leveraged to construct confidence estimations of retrieval performance. This idea of exploiting often unused information to improve retrieval is very related to the main body of this thesis's work. The confidence metrics have notably been integrated within the MTEB framework, as part of the MMTEB release which I am also a contributor of \sidecite{enevoldsen2025mmteb} (ICLR 2025). This thesis came at a time where open initiatives towards building multilingual LLMs were important and I participated in multiple such large projects following the success of CroissantLLM. Notably, with a team in Portugal, I worked on the EuroLLM series \sidecite{martins2024eurollmmultilinguallanguagemodels, martins2025eurollm} to train high performing open multilingual models with a special focus set on machine translation. I also contributed to an initiative to train improved multilingual bidirectional encoders which culminated in EuroBert \sidecite{boizard2025eurobert} (COLM 2025). With the EuroBert team, we further studied the differences between the training dynamics of encoder and decoder models in controlled setups, leading to \sidecite{gisserotboukhlef2025pretrainencodersmaskedlanguage}, under submission at ICLR 2026.

\noindent\textbf{ML Privacy Papers.} My early days of academic research were in the field of Computational Privacy. While I preferred focusing on LLMs and Information Retrieval during my thesis, privacy is a topic that stayed close to my heart. While working on CroissantLLM, I suggested to former colleagues that we could leverage this unique opportunity to also run LLM privacy research with huge LLM pretraining budgets. This led to a series of works, notably the Copyright Trap project \sidecite{meeus2024copyright} accepted at ICML 2024 in which we study model memorization of training data and draw conclusions into the risk of personal information leakage. A subsequent paper, \sidecite{meeus2024sok} provides a systematic  overview of membership inference attacks against LLM training data, organizing existing methods, clarifying threat models and assumptions, and highlighting open problems. It was awarded best paper award at IEEE SatML 2025. 

\noindent\textbf{Press \& Talks.}
Mention Press and Talks

\noindent\textbf{Ressource Releases.}
All models released as part of this thesis are openly released under open licenses.
